\begin{abstract}
This paper gives a technical account of Vibe Theory as it now stands: the mathematical
model, and a walk through thirty-six findings (P1 to P36) from a finite, reproducible
testbed. The model is monist: it treats reality as a single kind of entity, the
\emph{vibe}, a unit of experience carrying a ternary \emph{tone} and related to others
only by the \emph{note}, on a growing, random, hyperbolic, causal mesh, updated by a
local ternary signed-majority rule. There are no continuous weights and no hidden state,
and every continuous quantity of physics is an emergent, coarse-grained description of
the discrete substrate. The findings span the framework: geometry recovered sharply and
made the dominant phase of a measured transition, the law as a trilemma resolved by the
emergent-mesh Hamiltonian, a Lorentz-safe substrate, the four field spins with mass from
a Higgs and the gauge and graviton operators derived from the action, the Newtonian
limit and a conserving Einstein equation, the quantum pillars, and a cosmology with an
arrow, expansion, and a dark sector. Three predictions make contact with data, most
notably the everpresent cosmological constant matching the observed dark energy to order
of magnitude. Finally the model is run end-to-end, with the key emergent structures
coming off one instantiation. We are explicit about what is demonstrated, what is a first
rung on a deep problem, and what is open. This is a simulable model with a concrete build
target, not a finished empirical theory.
\medskip

\noindent\textbf{Status and caveats.} This is exploratory work, a feasibility study rather than a validated result. The aim is to see whether a model of this kind can be made to hang together at all, and so far it looks promising. Much validation and further work remain. Several of the experiments reported here are preliminary, and some may turn out to be inaccurate or improperly set up. Nothing here should be read as established or taken too rigorously at this stage. It is offered in the spirit of curiosity and invitation, a direction that looks worth exploring rather than a finished theory or a claim of discovery.
\end{abstract}
